\chapter{Methodology}
\label{ch:methodology}

\section{Overall Workflow}
\label{sec:overall_workflow}

The system follows a simulation-driven workflow. A user logs into the React dashboard, configures simulation parameters (seed value and strategy mode), and triggers a run. The backend creates an Orchestrator instance, which loads incident data, initializes four agents, and starts the message bus. Each agent processes its portion of the pipeline and passes structured messages to the next. When the final agent completes evaluation, run metrics are saved and the dashboard updates in real time.

Figure~\ref{fig:overall_workflow} illustrates this end-to-end flow.

\begin{figure}[H]
\centering
\fbox{\parbox{0.8\textwidth}{\centering\vspace{2cm}[Insert Overall Workflow Diagram]\vspace{2cm}}}
\caption{Overall system workflow from user input to result display}
\label{fig:overall_workflow}
\end{figure}

The workflow can be broken into five phases:

\begin{enumerate}[leftmargin=*]
    \item \textbf{Data Ingestion}: The Orchestrator checks whether any custom incidents have been reported through the dashboard. If custom incidents exist in the \texttt{custom\_incidents} table, those are loaded. Otherwise, the system reads from \texttt{data/input\_data.json}, which contains geospatially enriched synthetic disaster records generated by \texttt{generate\_data.py}.

    \item \textbf{Agent Pipeline Execution}: The Orchestrator publishes a \texttt{TRIGGER\_INCIDENT\_DETECTION} message to the bus. This triggers the four-stage pipeline: detection $\rightarrow$ allocation $\rightarrow$ rescue $\rightarrow$ evaluation.

    \item \textbf{State Machine Transitions}: The Orchestrator manages a \texttt{SystemState} enum with states \texttt{INIT}, \texttt{INCIDENT\_REPORTED}, \texttt{RESOURCE\_ALLOCATED}, \texttt{RESCUE\_IN\_PROGRESS}, \texttt{COMPLETED}, and \texttt{ERROR}. Each agent callback updates this state and logs the transition to the \texttt{run\_logs} table.

    \item \textbf{Real-Time Broadcasting}: Every state transition and inter-agent message is broadcast to connected WebSocket clients. The React dashboard receives these and updates the UI without page reloads.

    \item \textbf{Persistence and Export}: On completion, run metrics are saved to \texttt{run\_metrics}, and the dashboard enables CSV and PDF report downloads through dedicated API endpoints.
\end{enumerate}

\section{Agent Workflow}
\label{sec:agent_workflow}

The four agents execute sequentially, with each agent's output becoming the next agent's input. This is coordinated entirely through the message bus---agents do not call each other directly.

\begin{figure}[H]
\centering
\fbox{\parbox{0.8\textwidth}{\centering\vspace{2cm}[Insert Agent Pipeline Diagram]\vspace{2cm}}}
\caption{Sequential agent pipeline with message types}
\label{fig:agent_pipeline}
\end{figure}

\subsection{Incident Detection Agent}

This agent receives raw disaster data and produces an incident report for each record. In ``smart'' mode, it first attempts to use the trained ML model. The predictor loads three artifacts (\texttt{severity\_model.joblib}, \texttt{label\_encoder.joblib}, \texttt{feature\_scaler.joblib}) and returns a predicted severity label with a confidence score and per-class probabilities. If the model files are not present, the agent falls back to rule-based logic: high priority when severity is ``high'' or people affected exceeds 200, medium otherwise.

The output is an \texttt{IncidentReport} Pydantic model containing the affected area, latitude, longitude, priority, estimated risk, and ML prediction details. This output is wrapped in an \texttt{INCIDENT\_REPORTED} message and sent to the Resource Allocation Agent.

\subsection{Resource Allocation Agent}

Upon receiving incident data, this agent determines how many teams and which resources to assign. In smart mode, high-priority incidents receive 8 teams and a full resource set (ambulance, fire force, rescue boat, medical unit), while lower-priority incidents receive 3 teams with basic resources.

The agent then queries the SQLite \texttt{rescue\_teams} table for available teams, maps each team's coordinates and the incident's coordinates to the nearest city in the A* graph, and runs the solver to find the shortest path. The team with the lowest ETA is selected and its status is updated to ``busy'' in the database. The resulting allocation includes the route path, coordinates, distance in kilometers, and estimated arrival time.

\subsection{Rescue Coordination Agent}

This agent models the rescue operation. It reads team capacity from system configuration (18 per team in smart mode, 12 in baseline) and calculates how many people can be rescued given the number of assigned teams. A loss factor (0.08 for smart, 0.22 for baseline) introduces stochastic variability to simulate field-level uncertainty. The output is a \texttt{RescueStatus} containing the number of rescued people.

\subsection{Evaluation Agent}

The final agent computes efficiency metrics: lives saved as a percentage of people affected, response time, and resource utilization. It then calls the \texttt{DecisionSupportAgent}, which prompts Google Gemini with the full incident context. The Gemini response---or the local fallback---provides five structured fields: severity analysis, allocation analysis, evacuation recommendation, three concrete action items, and a two-sentence summary.

After processing all incidents, the Evaluation Agent returns overall simulation metrics including total incidents, total affected, total lives saved, overall efficiency percentage, and resource utilization status.

\section{Machine Learning Workflow}
\label{sec:ml_workflow}

\begin{figure}[H]
\centering
\fbox{\parbox{0.8\textwidth}{\centering\vspace{2cm}[Insert Machine Learning Pipeline Diagram]\vspace{2cm}}}
\caption{Machine learning training and inference pipeline}
\label{fig:ml_pipeline}
\end{figure}

The ML pipeline has two phases: offline training and online inference.

\subsection{Training Phase}

Training is triggered through the \texttt{POST /api/ml/train} endpoint, which runs the \texttt{train\_severity\_model.py} script in a background task. The script:

\begin{enumerate}[leftmargin=*]
    \item Generates 5000 synthetic training samples by randomly combining disaster types, Indian city characteristics, and correlated weather features. Flood incidents get higher rainfall, fire incidents get higher temperatures, and so on.

    \item Assigns severity labels using a domain-knowledge scoring function. Each disaster type has its own rules: floods score higher with heavy rainfall and dense populations, earthquakes score higher with poor infrastructure, fires score higher with high temperatures and wind speeds, and cyclones score higher with extreme winds.

    \item Extracts features through the \texttt{feature\_engineering.py} module. This produces a 13-dimensional feature vector per sample: 4 one-hot columns for disaster type, plus 9 numeric features (people affected, temperature, humidity, wind speed, rainfall, population density, infrastructure score, hour of day, month).

    \item Splits the data 80/20 with stratification on the severity label.

    \item Trains a \texttt{RandomForestClassifier} using \texttt{GridSearchCV} with 3-fold cross-validation. The grid searches over \texttt{n\_estimators} $\in \{100, 200\}$, \texttt{max\_depth} $\in \{10, 20, \text{None}\}$, and \texttt{min\_samples\_split} $\in \{2, 5\}$.

    \item Saves the best model, label encoder, and feature scaler as \texttt{.joblib} files, along with a \texttt{training\_report.json} containing accuracy, F1 score, feature importances, and best hyperparameters.
\end{enumerate}

\subsection{Inference Phase}

At runtime, the \texttt{SeverityPredictor} singleton loads the saved artifacts. For each incident, it runs \texttt{extract\_features()} to build the feature vector, applies the scaler, calls \texttt{predict\_proba()} and \texttt{predict()}, and returns the predicted label with confidence and per-class probabilities. This prediction is attached to the \texttt{IncidentReport} as \texttt{ml\_predicted\_severity}, \texttt{ml\_confidence}, and \texttt{ml\_probabilities}.

\section{Database Workflow}
\label{sec:database_workflow}

SQLite serves as the sole persistence layer. The \texttt{init\_db()} function runs at application startup and creates ten tables with \texttt{CREATE TABLE IF NOT EXISTS} statements, ensuring idempotent initialization. A separate migration function (\texttt{\_migrate\_existing\_tables}) adds new columns to existing tables without data loss.

\begin{figure}[H]
\centering
\fbox{\parbox{0.8\textwidth}{\centering\vspace{2cm}[Insert Database Interaction Diagram]\vspace{2cm}}}
\caption{Database read/write interactions across the simulation lifecycle}
\label{fig:db_workflow}
\end{figure}

During a simulation run, the database is written to at multiple points:

\begin{itemize}[leftmargin=*]
    \item \texttt{create\_run()} inserts a row into \texttt{simulation\_runs} when the Orchestrator initializes.
    \item \texttt{log\_transition()} inserts into \texttt{run\_logs} at every state machine transition.
    \item \texttt{save\_incidents()} inserts one row per incident into the \texttt{incidents} table.
    \item \texttt{save\_agent\_log()} records each agent's output as a JSON string in \texttt{agent\_logs}.
    \item \texttt{save\_agent\_message()} logs every message published on the bus to \texttt{agent\_messages}.
    \item \texttt{update\_incident\_routing()} updates incident rows with assigned team ID, route path, distance, and ETA.
    \item \texttt{save\_run\_metrics()} inserts aggregate metrics on completion.
    \item \texttt{complete\_run()} updates the run status to ``COMPLETED'' or ``ERROR'' with a timestamp.
\end{itemize}

\section{Communication Workflow}
\label{sec:communication_workflow}

\begin{figure}[H]
\centering
\fbox{\parbox{0.8\textwidth}{\centering\vspace{2cm}[Insert Message Bus Workflow Diagram]\vspace{2cm}}}
\caption{Message bus publish-subscribe-route cycle}
\label{fig:message_bus_workflow}
\end{figure}

The message bus is the backbone of inter-agent communication. It is implemented in \texttt{message\_bus.py} as a thread-safe, priority-queued publish-subscribe system.

\subsection{Message Structure}

Each message contains a unique ID (UUID), the run ID, sender name, receiver name, message type, payload dictionary, priority level (1--5, where 1 is critical), timestamp, and retry count. Messages implement \texttt{\_\_lt\_\_} for priority queue ordering: lower priority numbers are dequeued first, with timestamp as a tiebreaker.

\subsection{Publish-Route-Handle Cycle}

When an agent or the Orchestrator publishes a message, the bus:
\begin{enumerate}
    \item Assigns a UUID and constructs a \texttt{Message} object.
    \item Persists the message to the \texttt{agent\_messages} table.
    \item Broadcasts the message to WebSocket clients.
    \item Enqueues the message in a \texttt{PriorityQueue}.
\end{enumerate}

The \texttt{process\_next()} method dequeues the highest-priority message, looks up the receiver's registered handler, calls it, and publishes any reply messages generated by the handler. If the receiver is \texttt{BROADCAST}, the message is sent to all registered agents except the sender.

\subsection{Retry and Health Monitoring}

If a handler raises an exception, the bus increments the message's retry count and re-enqueues it after a configurable delay (default 10ms). After three failed attempts, the message is dropped and logged. Each agent has an \texttt{AgentHealth} record tracking messages processed, messages failed, cumulative latency, and a status flag (healthy, degraded, unhealthy).

\section{Project Workflow}
\label{sec:project_workflow}

The project was developed iteratively over several phases:

\begin{enumerate}[leftmargin=*]
    \item \textbf{Phase 1 -- Core Agent Framework}: Built the four agents with rule-based logic, the Orchestrator with direct method calls, and the SQLite schema with basic tables. Verified the pipeline end-to-end.

    \item \textbf{Phase 2 -- Message Bus Integration}: Replaced direct agent calls with the centralized message bus. Added priority queuing, retry logic, and health tracking. Updated the Orchestrator to use wrapped callbacks.

    \item \textbf{Phase 3 -- Machine Learning}: Developed the synthetic data generator, feature engineering module, and Random Forest training pipeline. Integrated the predictor into the Incident Detection Agent with graceful fallback.

    \item \textbf{Phase 4 -- Route Optimization}: Built the A* solver with Haversine heuristics and a graph of Indian highway connections. Integrated it into the Resource Allocation Agent to select the closest available team.

    \item \textbf{Phase 5 -- Gemini AI Integration}: Created the \texttt{DecisionSupportAgent} with modular LLM provider architecture. Added caching, JSON response parsing, and the local fallback generator.

    \item \textbf{Phase 6 -- Frontend and Real-Time Updates}: Built the React dashboard with WebSocket integration, Leaflet map, SVG charts, and report download functionality. Styled the interface with a dark theme and animated elements.
\end{enumerate}

\section{Chapter Summary}
\label{sec:ch3_summary}

This chapter described the methodology in terms of five interleaved workflows: the overall simulation flow, the agent interaction model, the ML training and inference pipeline, the database persistence pattern, and the message bus communication cycle. The project was developed in six iterative phases, starting from core agents and incrementally adding ML, routing, AI integration, and the frontend dashboard.
